\subsection{Out-of-Design Scenario Stress Test}
\label{sec:out-of-design-stress}

The static frontier evaluation in Section~\ref{sec:results} demonstrates that BDT anchors unconstrained LLM agents toward DCE-implied behavioral consistency within the original 64-state attribute grid. However, a natural question arises: does BDT's anchoring come at the cost of limiting the LLM's ability to respond to richer policy descriptions that fall outside the DCE design space?

To address this, we construct a stress test using 12 scenarios that vary only on dimensions \emph{not} represented in the original DCE. Each scenario holds the DCE-measured attributes fixed (wait time 2 or 6 months, vaccine efficacy 70\%, side effect risk low) while varying an auxiliary policy description. The six auxiliary conditions are: (a) neutral standard appointment, (b) online pre-registration required, (c) walk-in access available, (d) appointment site far from home, (e) appointment reminder with flexible scheduling, and (f) uncertain queue length with administrative paperwork.

We obtain $P_{\mathrm{static}}$ from the mixed-logit frontier using only the DCE-measured dimensions ($\beta_{\mathrm{wait}} = -0.267$, $\beta_{\mathrm{eff}} = 1.459$, $\beta_{\mathrm{se}} = -0.102$). All 12 scenarios produce $P_{\mathrm{static}} = 0.619$ at 2 months wait and $P_{\mathrm{static}} = 0.358$ at 6 months, since the auxiliary descriptions are not represented in the DCE design. We then obtain unconstrained LLM probabilities $p_{\mathrm{LLM}}$ by prompting the model with the full scenario text (10 repetitions), and compute BDT-anchored probabilities as $\pi_{\mathrm{BDT}} = \lambda p_{\mathrm{LLM}} + (1 - \lambda) P_{\mathrm{static}}$ with $\lambda = 0.25$.

\paragraph{Results.}
Table~\ref{tab:out-of-design} reports the stress-test metrics. The unconstrained LLM exhibits an MVR-Wait (monotonicity violation rate) of 0.33 across the six wait-time pairs, meaning that in two out of six paired comparisons, the LLM assigned a \emph{higher} acceptance probability to the 6-month wait than to the 2-month wait. The mean absolute deviation (MAD) from $P_{\mathrm{static}}$ is 0.290, and the Spearman rank correlation with $P_{\mathrm{static}}$ is near zero ($\rho = 0.12$, $p = 0.71$). These findings confirm that unconstrained LLM agents do not reliably reflect DCE-measured attribute trade-offs when presented with richer textual scenarios.

After BDT anchoring, MVR-Wait is fully resolved (no monotonicity violations across all six pairs). MAD from $P_{\mathrm{static}}$ decreases substantially to 0.072 (a 75\% reduction), and the Spearman correlation increases correspondingly. This demonstrates that BDT preserves the DCE-implied behavioral discipline---specifically the ordering over wait time---while allowing the LLM to adjust its probability mass in response to the auxiliary policy descriptions.

\paragraph{Interpretation.}
This stress test is not intended as external predictive validation. Rather, it evaluates whether BDT can simultaneously satisfy two competing objectives: (1) maintaining consistency with DCE-measured behavioral regularities, and (2) preserving the flexibility of LLM-based scenario representation for dimensions not directly encoded in the DCE design. The results indicate that BDT achieves both: the anchoring term ensures behavioral discipline on the DCE-measured dimensions, while the LLM term ($\lambda = 0.25$) allows agents to modulate their responses based on richer policy descriptions. The auxiliary descriptions themselves show plausible directional effects---walk-in access and flexible scheduling increase acceptance, while administrative paperwork reduces it---suggesting that the LLM component captures meaningful policy-relevant variation outside the DCE design space.
